\section{第\ref{sec:gaba}章　\nt{GLUT}与\nt{GABA}协同工作}

本章的逻辑命题和动力学命题都在本章中由线性分析和欧姆定律推出；实验表明，它们所依据的电路——带延迟的否决、紧随兴奋之后的前馈抑制、由抑制实现的稳定、\nt{GLUT}--\nt{GABA}环路的节律——确实在大脑中工作。

{\footnotesize
\setlength{\tabcolsep}{3pt}\renewcommand{\arraystretch}{1.15}
\begin{longtable}{>{\raggedright}p{0.5cm}>{\raggedright}p{4.0cm}>{\raggedright}p{5.6cm}>{\raggedright}p{2.3cm}>{\raggedright\arraybackslash}p{1.7cm}}
\toprule
编号 & 命题 & 实验与测量内容 & 来源 & 状态 \\
\midrule\endfirsthead
\toprule
编号 & 命题 & 实验与测量内容 & 来源 & 状态 \\
\midrule\endhead
1 & 皮层中五分之一到四分之一的神经元是\nt{GABA}神经元 & 猴皮层 $10$ 个区的GABA免疫染色：$8$ 个区中约占神经元的 $25\,\%$，初级视皮层中不到 $20\,\%$。皮层抑制性神经元多样性综述 & \cite{Hendry1987,Markram2004} & 已测量 \\
2 & 抑制的作用在很大程度上取决于落点：树突、胞体、脉冲产生处、别的导线的末梢 & 综述：皮层\nt{GABA}神经元的各个类别按其在接收者上的靶点区分。同时记录锥体神经元的胞体和树突：前馈抑制在胞体上比在树突上强得多。脊髓中作用于别的导线末梢的抑制减少一条输入线路的递质释放（测量综述） & \cite{Markram2004,Pouille2001,Rudomin1999} & 已测量 \\
3 & 经中介改变符号的代价是一个延迟，约一毫秒；紧随兴奋而来的抑制收窄符合窗口 & 大鼠海马脑片中的锥体神经元：经一个中介的抑制以很短的延迟紧随直接兴奋到达，两个兴奋性输入能总和到阈值的窗口小于 $2\ms$。在这种抑制较弱的树突上，窗口较宽 & \cite{Pouille2001} & 已测量 \\
4 & 否决 $a\wedge\bar b$~\eqref{eq:veto} 加上或和常数1是函数完备的（命题~\ref{prop:gabacomplete}） & 本章的证明。经典结果：带抑制性输入的阈值元件网络可以实现任意逻辑表达式。这是关于模型的命题，没有实验 & \cite{McCullochPitts1943} & 理论 \\
5 & 带延迟的否决给出运动方向检测器，最优速度 $v^*=\Delta x/d$~\eqref{eq:vstar} & 兔视网膜：方向选择性由抑制产生，抑制在“零”方向运动时抵消兴奋。分别记录方向选择细胞的兴奋性和抑制性输入，表明星爆无长突细胞（\nt{GABA}）对它们的抑制是不对称的，只来自朝向“零”方向一侧的突起。$v^*$ 的公式是本章的推导 & \cite{Barlow1965,Fried2002} & 已测量 \\
6 & 分流抑制对电压做除法而非减法~\eqref{eq:shunt}（图~\ref{fig:shunt}） & 三个电导构成的分压器的欧姆定律，氯的平衡电位在静息附近；针对不发放脉冲的神经元 & 本章推导 & 理论 \\
7 & 分流对脉冲频率的作用是减法，除法来自抑制加上平衡的背景噪声 & 模型：发放中的神经元流过分流通道的电流几乎恒定，对频率的效应是减法。实验：用动态钳向大鼠体感皮层锥体神经元（脑片）注入兴奋性和抑制性背景电导；两者同时平衡地增加，会除以响应斜率，而频率没有明显平移。综述：除以总活动的运算见于大脑许多部位 & \cite{HoltKoch1997,Chance2002,Carandini2012} & 已测量 \\
8 & $\beta>1$ 时相互抑制选出唯一赢家（命题~\ref{prop:wta}，\eqref{eq:wta}） & 特征值 $-(1\pm\beta)/\tau$ 是本章的推导。$\beta=0.5$ 时的频率 $0.73$ 和 $0.53$ 是不动点 $r_{1,2}=(I_{1,2}-\beta I_{2,1})/(1-\beta^2)$；\texttt{models/} 中没有对应脚本。本章未引用直接实验 & 本章推导 & 理论 \\
9 & 经\nt{GABA}用信号复位存储；两群相互抑制的神经元构成锁存器 & 由图~\ref{fig:bistable} 和命题~\ref{prop:wta} 推出。没有实验 & 本章推导 & 理论 \\
10 & 若抑制足够强且足够快，\nt{GLUT}--\nt{GABA}环路的平衡点稳定（命题~\ref{prop:ei}） & 双群体模型~\eqref{eq:ei}；条件 (i) 和 (ii) 是其矩阵的行列式和迹，在本章中推导 & \cite{WilsonCowan1972} & 理论 \\
11 & $w_{EE}=1.5$、$w_{EI}w_{IE}=2$ 时，快速抑制（$\tau_I=2\ms$）维持 $E^*=0.67h$，慢速抑制（$30\ms$）激起振荡（图~\ref{fig:ei}） & \eqref{eq:ei} 的数值解，脚本 \texttt{figures/make\_ei.py} & 计算 & 模型 \\
12 & 皮层工作在抑制稳定状态；其特征是推一下\nt{GABA}神经元，它们的频率反而下降~\eqref{eq:isn} & 理论预言了这种反常响应。实验：猫初级视皮层的细胞内记录——当响应被周围刺激抑制时，抑制性电导不增反降，与兴奋性电导一同下降。在小鼠视皮层、体感皮层和运动皮层中光遗传学兴奋PV神经元，会降低抑制性神经元的频率，包括无刺激时和浅麻醉下。图~\ref{fig:ei} 数值下的系数 $-1/3$ 是本章的推导 & \cite{Tsodyks1997isn,Ozeki2009,Sanzeni2020} & 已测量 \\
13 & “\nt{GLUT}兴奋\nt{GABA}，\nt{GABA}抑制\nt{GLUT}”环路产生周期为 $12$--$50\ms$ 的快速节律 & 在小鼠体感皮层以 $8$--$200$\,Hz光遗传学刺激快速\nt{GABA}神经元，选择性地增强 $\gamma$ 节律（$20$--$80$\,Hz，周期 $12$--$50\ms$）；刺激\nt{GLUT}神经元只增强较慢的节律 & \cite{Cardin2009,Buzsaki2012} & 已测量 \\
\bottomrule
\end{longtable}}
